%%%%%%%%%%%%%%%%%%%%%%%%%%%%%%%%%%%%%%%%%%%%%%%%%%%%%%%%%%%%%%%%%%%%%%%%%%%%%%%%
%
% Plantilla para libro de texto de matemáticas.
%
% Esta plantilla ha sido desarrollada desde cero para el proyecto apuntesDGIIM
% de LibreIM.
%
%	Copyright (C) 2019 LibreIM
%
% This program is free software: you can redistribute it and/or modify
% it under the terms of the GNU Affero General Public License as
% published by the Free Software Foundation, either version 3 of the
% License, or (at your option) any later version.
% This program is distributed in the hope that it will be useful,
% but WITHOUT ANY WARRANTY; without even the implied warranty of
% MERCHANTABILITY or FITNESS FOR A PARTICULAR PURPOSE.  See the
% GNU Affero General Public License for more details.
% You should have received a copy of the GNU Affero General Public License
% along with this program.  If not, see <https://www.gnu.org/licenses/>.
%
%%%%%%%%%%%%%%%%%%%%%%%%%%%%%%%%%%%%%%%%%%%%%%%%%%%%%%%%%%%%%%%%%%%%%%%%%%%%%%%%

% ------------------------------------------------------------------------------
% CONFIGURACIÓN BÁSICA DEL DOCUMENTO
% ------------------------------------------------------------------------------

\documentclass[fontsize=12pt]{scrartcl}
\usepackage[english]{babel}
\usepackage{graphicx}

\usepackage{templatetools}

% ------------------------------------------------------------------------------
% CONFIGURACIÓN DE ASIGNATURA
% ------------------------------------------------------------------------------

\usepackage{config}

% ------------------------------------------------------------------------------
% TIPOGRAFÍA
% ------------------------------------------------------------------------------

% Tipografía personalizada
\usepackage[bitstream-charter]{mathdesign}
\usepackage[scale=0.88, type1]{FiraSans}
\usepackage[scale=0.88, type1]{FiraMono}
\usepackage[T1]{fontenc}

% Microtype
\usepackage[activate={true, nocompatibility}, final, tracking=true,factor=1100, stretch=10, shrink=10]{microtype}
\SetTracking{encoding={*}, shape=sc}{0}

% ------------------------------------------------------------------------------
% DISEÑO DE PÁGINA
% ------------------------------------------------------------------------------

% Márgenes
\usepackage[bottom=3.125cm, top=2.5cm, left=2.5cm, right=4.5cm, marginparwidth=70pt]{geometry}

% Párrafos
\linespread{1.1}
\setlength{\parindent}{12pt}
\setlength{\parskip}{0pt}

% Listas
\setlist{noitemsep}

% Cabeceras
\usepackage[automark]{scrlayer-scrpage}
\clearpairofpagestyles
\chead{\leftmark}
\cfoot*{\pagemark}

% Enlaces
\usepackage{hyperref}
\hypersetup{
	colorlinks=true,
	linkcolor=500,
}

% No mostrar cabecera en páginas de sección
\usepackage{xpatch}
\xpretocmd{\section}{\vspace*{1cm}\thispagestyle{plain}}{}{}

% Páginas de color para portada
\usepackage{pagecolor}
\usepackage{afterpage}

% ------------------------------------------------------------------------------
% ENTORNOS PERSONALIZADOS
% ------------------------------------------------------------------------------

\usepackage[theorems, skins, breakable]{tcolorbox}

\tcolorboxenvironment{nth}{
	blanker,
	breakable,
	left=12pt,
	before skip=12pt,
	after skip=12pt,
	borderline west={2pt}{0pt}{500},
	before upper={\parindent 12pt},
}

\tcolorboxenvironment{nprop}{
	blanker,
	breakable,
	left=12pt,
	before skip=12pt,
	after skip=12pt,
	borderline west={2pt}{0pt}{300},
	before upper={\parindent 12pt},
}

\tcolorboxenvironment{ncor}{
	blanker,
	breakable,
	left=12pt,
	before skip=12pt,
	after skip=12pt,
	borderline west={2pt}{0pt}{300},
	before upper={\parindent 12pt},
}

\tcolorboxenvironment{ndef}{
	skin=enhancedmiddle jigsaw,
	frame hidden,
	colback=50,
	breakable = true,
	break at = -6pt,
	top = 4pt,       % Estos márgenes están un poco a ojo
	bottom = 4pt,
	left= 8pt,
	right = 8pt,
	before skip=8pt, % Normalmente dejamos 12pt, pero
	after skip=8pt,  % aquí tenemos espacio adicional por el fondo
	no borderline,
	borderline west={2pt}{0pt}{500},
	borderline east={2pt}{0pt}{50},
	before upper={\parindent 12pt},
}

\tcolorboxenvironment{ejer}{
	skin=enhancedmiddle jigsaw,
	frame hidden,
	colback=50,
	breakable = true,
	break at = -6pt,
	top = 4pt,       % Estos márgenes están un poco a ojo
	bottom = 4pt,
	left= 8pt,
	right = 8pt,
	before skip=8pt, % Normalmente dejamos 12pt, pero
	after skip=8pt,  % aquí tenemos espacio adicional por el fondo
	no borderline,
	borderline west={2pt}{0pt}{500},
	borderline east={2pt}{0pt}{50},
	before upper={\parindent 12pt},
}

\tcolorboxenvironment{remark}{
	blanker,
	breakable,
	left=12pt,
	before skip=12pt,
	after skip=12pt,
	borderline west={2pt}{0pt}{300},
	before upper={\parindent 12pt},
}


% ------------------------------------------------------------------------------
% LISTINGS
% ------------------------------------------------------------------------------

\IfPackageLoaded{listings}{
	\renewcommand{\lstlistingname}{Listado}

	\lstset{
		frame=lines,
		rulecolor=\color{black},
		framerule=1pt,
		numbers=left,
		belowcaptionskip=1\baselineskip,
		basicstyle=\ttfamily\color{black},
		keywordstyle=\bfseries\color{700},
		commentstyle=\color{300},
		stringstyle=\color{500},
		numberstyle=\color{black},
		breaklines=true,
		showstringspaces=false,
		tabsize=2,
	}
}


\usepackage[utf8]{inputenc}
\usepackage{pgfplots}
\DeclareUnicodeCharacter{2212}{−}
\usepgfplotslibrary{groupplots,dateplot}
\usetikzlibrary{patterns,shapes.arrows}
\pgfplotsset{compat=newest}
\usepackage{float}

%%%%%%%%%%%%%%%%%%%%%%%%%%%%%%%%%%%%%%%%%%%%%%%%%%%%%%%%%%%%%%%%%%%%%%%%%%%%%%%%
% ------------------------------------------------------------------------------
% COMIENZO DEL DOCUMENTO
% ------------------------------------------------------------------------------
%%%%%%%%%%%%%%%%%%%%%%%%%%%%%%%%%%%%%%%%%%%%%%%%%%%%%%%%%%%%%%%%%%%%%%%%%%%%%%%%

\begin{document}

% ---------------------------------------------------------------------------
% PORTADA EXTERIOR
% ---------------------------------------------------------------------------

\newpagecolor{500}
\begin{titlepage}
	\noindent
	\setlength\fboxsep{0cm}
	\parbox[t]{\textwidth}{
			\raggedright
			\fontsize{50pt}{50pt}\selectfont\sffamily\color{white}{
			  \textbf{\subj}
      }
	}

	\vfill

	\noindent
	\parbox[b]{\textwidth}{
		\raggedright
		\sffamily\large\color{white}
		{\Large \autor }\\[4pt]
		\grado\\
		\universidad\\[4pt]
		\texttt{\enlaceweb}
	}

\end{titlepage}
\restorepagecolor

% ---------------------------------------------------------------------------
% PÁGINA DE LICENCIA
% ---------------------------------------------------------------------------

\thispagestyle{empty}
\null
\vfill

\noindent
\parbox[b]{0.7\textwidth}{
  \includegraphics{by-nc-sa.pdf}\\[4pt]
  \raggedright
  \sffamily
  {\large This book is distributed under a license CC BY-NC-SA 4.0.}\\[4pt]
  \texttt{creativecommons.org/licenses/by-nc-sa/4.0/}
}

\newpage
% ---------------------------------------------------------------------------
% ÍNDICE
% ---------------------------------------------------------------------------

\thispagestyle{empty}
{\hypersetup{linkcolor=black}\tableofcontents}
\newpage

% ---------------------------------------------------------------------------
% CONTENIDO
% ---------------------------------------------------------------------------

\part{Introduction}
\section{Time Series}

A time series, also known as discrete time signal, is a sequence of observations taken periodically in time. We can use time series to perform many tasks such as predictions of future values, behaviour analysis or information extraction. Examples of time series are audio signals, industrial instrument measures or diary finantial activity.

A system can be determined comparing the input and the output. We call the system a filter if it is linear and time invariant. Considering the dynamic system as a black box, we can estimate the transference function or the impulse response to taht filter.

We can also consider \textbf{multivariate} time series, where some values of the time series have an influence on the other values in different or the same time instant. We can \textbf{classify} the time series in two wide types:

\begin{itemize}
  \item Determinist: based in dynamic systems, they exploit the phisics of the generation algorithm of the time series.
        \item Stochastic: where the series are realizations of a stochastic process, which can be modelated.
\end{itemize}

In this subject, we will focus on stochastic models.

\newpage

\end{document}
